% Internal route lemma only. Not promoted or counted as original A.5 movement.
\subsection*{Double mapping cylinders and coherent definable approximation}

Let $R$ be a real closed field with an o-minimal expansion. Let
$F\subset R^m$ and $G\subset R^n$ be closed bounded definable sets,
and let $q:F\to G$ be a continuous definable surjection. Define
\[
 Z_q=\{(s,sx,q(x)):s\ne0,\ x\in F\}
         \cup\{(0,0,g):g\in G\}\subset R\times R^m\times R^n,
 \qquad p(s,z,g)=s.
 \tag{MC.1}
\]
This is a definable realization of the double mapping cylinder of $q$.
The map $p$ is continuous and definably proper. Its fibres for $s\ne0$
are copies of $F$, and its zero fibre is a copy of $G$.

To check closedness and properness directly, consider
$\Phi:R\times F\to R\times R^m\times R^n$ given by
$\Phi(s,x)=(s,sx,q(x))$. Its image is $Z_q$. A closed bounded set
in the target has bounded first coordinate, and its inverse image under
$\Phi$ is a closed subset of a closed bounded product in $R\times F$.
Thus $\Phi$ is definably proper, so its image is closed. If $K\subset R$
is closed bounded and definable, then
$p^{-1}(K)=\Phi(K\times F)$ is closed bounded, proving the claim.

\paragraph{Precise single-map criterion.}
The following two conditions are equivalent.
\begin{enumerate}
\item The map $p$ is definably locally trivial at $0$.
\item There exist $\epsilon>0$ and a continuous definable map
\[
 h:(0,\epsilon)\times G\longrightarrow F
\]
such that each $h_s:G\to F$ is a homeomorphism and
\[
 \lim_{s\to0^+}\sup_{g\in G}\lVert q(h_s(g))-g\rVert=0.
 \tag{MC.2}
\]
\end{enumerate}
We prove the criterion, including the uniform convergence requirement.

Suppose first that a product chart
$H:(-\epsilon,\epsilon)\times G\to p^{-1}((-
\epsilon,\epsilon))$ is given. Its restriction to the zero fibre is
$(0,g)\mapsto(0,0,k(g))$ for some definable homeomorphism $k:G\to G$.
Precomposing the chart with $(s,g)\mapsto(s,k^{-1}(g))$ makes this
restriction the identity. For $s>0$ write
\[
 H(s,g)=(s,s h_s(g),q(h_s(g))).
\]
The maps $h_s$ are homeomorphisms because the coordinate map
$x\mapsto(s,sx,q(x))$ identifies the $s$-fibre with $F$.
Their joint map is definable and continuous for $s>0$, by recovering
$h_s(g)$ as the middle coordinate divided by $s$.
The last coordinate of the normalized chart extends continuously to
$(0,g)$ with value $g$. Definable compactness of $G$ makes this
convergence uniform: otherwise definable choice supplies $g(s)$ with
error bounded away from zero, and after shortening the interval each
coordinate of the bounded definable curve $g(s)$ has a limit
$g_0\in G$, contradicting joint continuity at $(0,g_0)$.
This proves (MC.2).

Conversely, given $h$ as in (2), define
\[
 H(s,g)=
 \begin{cases}
   (s,s h_{|s|}(g),q(h_{|s|}(g))),&s\ne0,\\
   (0,0,g),&s=0.
 \end{cases}
 \tag{MC.3}
\]
It is definable and bijective onto $p^{-1}((-
\epsilon,\epsilon))$. It is continuous away from zero. At zero,
the middle coordinate tends uniformly to zero because $F$ is bounded,
and the last coordinate tends uniformly to $g$ by (MC.2). Thus $H$
is continuous.

Its inverse is continuous away from zero as well. Indeed, for each
closed subinterval $J\subset(0,\epsilon)$ the map
$(s,g)\mapsto(s,h_s(g))$ is a continuous definable bijection from
the closed bounded set $J\times G$ onto $J\times F$, and so it is
a homeomorphism. These intervals cover the nonzero parameter locus,
with the negative side obtained using $|s|$. Finally, if
$H(s,g)=(s,z,w)$ with $s\ne0$, then
\[
 \lVert g-w\rVert
       \le\sup_{y\in G}\lVert y-q(h_{|s|}(y))\rVert.
\]
The right side tends to zero. Hence $(s,z,w)\to(0,0,g_0)$ forces
$g\to g_0$, proving inverse continuity at every zero-fibre point.
This establishes the product chart.

\paragraph{What remains missing.}
The condition in (2) is a \emph{jointly definable} family of
homeomorphisms. It must not be replaced by the weaker assertion that
for every $\delta>0$ there exists some definable homeomorphism with
error at most $\delta$. Such individually existing homeomorphisms may
have unbounded formula complexity, and definable choice cannot be
applied before a definable parameter space of candidates is supplied.
The criterion therefore exhibits the coherent approximation issue; it
does not solve it.

If $q$ is replaced by a continuous definable family
$q_a:F\to G$ over a definable parameter set $A$, the same realization
\[
 Z=\{(a,s,sx,q_a(x)):s\ne0,\ x\in F\}
       \cup\{(a,0,0,g):a\in A,\ g\in G\}
 \tag{MC.4}
\]
is definably proper over $A\times R$. A local product at $(a_0,0)$
gives, after normalizing its central family of homeomorphisms,
homeomorphisms $h_{a,s}:G\to F$ which are jointly definable and
continuous, with $q_a h_{a,s}$ extending to the identity at $s=0$.
Conversely such a joint family gives a local product by (MC.3), now
with the additional parameter $a$. Uniformity in $g$ near $(a_0,0)$
follows by the same definable-choice argument: along a definable curve
of violations, the bounded $G$-coordinate has coordinatewise limits
in $G$, and the parameter coordinate already tends to $a_0$. Joint
continuity at the limiting point gives the contradiction. No compact
parameter neighbourhood or ordinary subsequence is required.
This is a reduction of a special class of A.5 maps to a coherent
approximation problem, not a proof of that problem or a reduction of
the whole A.5 question to it.

No external theorem beyond elementary definable compactness is used in
this route lemma. No finite combinatorial search is asserted or needed.

